% ----------------------------- RESUMEN -----------------------------
\clearpage
\blank
\caratula{RESUMEN}
\begingroup\dobleesp
\vspace{6pt}
Cuando cambian el modelo, el prompt o las herramientas de un agente conversacional basado en LLM, no hay una forma establecida de saber si la nueva versión se comporta peor que la anterior: sus respuestas varían entre ejecuciones y una conversación de varios turnos no se repite palabra por palabra. Este trabajo presenta AgentProof, una herramienta de código abierto que induce de forma automática un escenario a partir de cada conversación que el agente ya sostuvo, lo reejecuta varias veces contra la versión nueva con un usuario simulado, juzga cada repetición a ciegas con un evaluador automático y declara una regresión solo con evidencia estadística, corregida por comparaciones múltiples. Su capacidad de detección se midió con un protocolo pre-registrado sobre un agente de atención al cliente, veinte conversaciones semilla elicitadas y seis degradaciones introducidas como mutaciones. En 1200 conversaciones reejecutadas, una política de reembolsos corrupta se detectó en todos los escenarios donde se manifestó, con el agente construido en dos frameworks, y el control negativo no declaró ninguna regresión. En cambio, dos degradaciones cuyo efecto se reparte entre escenarios bajaron la aprobación de la suite entre 25 y 30 puntos sin declararse, por la resolución del criterio. Tres anotadores externos validaron el evaluador, que coincide con su mayoría en 34 de 40 conversaciones ($\kappa = 0.70$), y la fidelidad de los escenarios inducidos. El agente se construyó además con smolagents y CrewAI sin modificar el núcleo, y la misma mutación del límite de pasos se manifestó distinto según el framework. Los resultados indican que la reejecución con usuarios simulados detecta degradaciones sistemáticas sin alarmas cuando el agente no cambia, y que ante degradaciones repartidas el reporte debe leerse junto con la caída de la suite.
\par\vspace{6pt}
\textbf{Palabras clave:} agentes basados en LLM, detección de regresiones, usuarios simulados, LLM como evaluador, pruebas de mutación, inducción de escenarios, conversaciones multi-turno, pre-registro.
\par
\endgroup
